\section{Method}
\label{sec:method}

\subsection{Task}
\label{sec:method:task}

A model is given a user request and a set of tool schemas in its own chat format, and generates greedily. The
generated span either contains a tool call or does not. We label the call by comparing it against the benchmark's
ground truth after normalising both sides to a canonical name and argument dictionary, and we record a
\emph{failure mode} rather than a bit: \texttt{valid}, \texttt{wrong\_name}, \texttt{wrong\_arg\_values},
\texttt{missing\_args}, \texttt{missing\_calls}, \texttt{no\_call} (prose instead of a call),
\texttt{unparseable\_call}, and, on the benchmark category where no tool applies, \texttt{valid\_nocall} and
\texttt{over\_trigger}. The binary label is \texttt{valid} and \texttt{valid\_nocall} against the rest.

The modes are not bookkeeping. Two of them, \texttt{no\_call} and \texttt{unparseable\_call}, are format
failures that a length feature can detect without any access to internals, so every result is reported both on
the full set and on the \emph{semantic} subset that excludes them, where the detector must separate a correct
call from a well-formed but wrong one. On the benchmark that mixes a category with no applicable tool, we report
a third column restricted to the \emph{call-expected} items, because on the remaining items the label reduces to
whether a call was emitted at all.

\subsection{Access tiers}
\label{sec:method:tiers}

Deployments differ in what they can read from a model, and the difference is not cosmetic. Token
log-probabilities are returned by every serving API. Residual-stream states are a forward hook away in any
stack that runs the model's own code, but are not returned by hosted APIs. Attention probabilities are the
most awkward of the three: fused attention kernels never materialise them, so reading them requires the eager
attention path or a second pass, which is what we do (Appendix~\ref{app:protocol}). The attention tier is
therefore not the cheapest tier to reach; its interest is that it needs no activations, so a judge built on it
can be trained and audited from attention maps alone, and that its features have a graph-theoretic reading
that the theory of Section~\ref{sec:theory} can address. We group detectors by what they consume and report a
frontier over tiers rather than a single ranking, and we measure what each costs at the moment the call is
produced (Appendix~\ref{app:cost}).

\begin{description}[topsep=2pt,itemsep=1pt,leftmargin=1.2em]
\item[Logits.] Mean token log-probability of the generated span, with the sign fixed on training folds only.
\item[Attention.] Functions of the attention weight matrices alone (Section~\ref{sec:method:attn}).
\item[Residual.] Functions of the residual-stream states, that is, the block outputs that
\texttt{output\_hidden\_states} returns (Section~\ref{sec:method:resid}).
\end{description}

We keep the tiers strictly separate, which requires care. A family of published ``attention'' diagnostics
projects residual-stream vectors onto the attention eigenbasis, giving Dirichlet energy, smoothness, spectral
entropy and high-frequency energy ratio \citep{noel2025gsp}. Those quantities require both tensors and cannot be
computed under attention-only access. In prior work built on that library, four of five nominally spectral
features were of this hybrid kind, which inflates an attention-only claim with residual-stream information. Appendix~\ref{app:provenance} tags every feature by the tensors it consumes.

\subsection{Attention-only detectors}
\label{sec:method:attn}

All of the following are computed on the subgraph induced on the generated call, using
Eq.~\ref{eq:adj} and the normalised Laplacian.

\paragraph{Per-head spectral profile (ours).}
For each layer $\ell$ and head $h$ we take one eigendecomposition of $\mathcal{L}_{\ell h}$ and read five
eigenvalue-only quantities: the Fiedler value $\lamtwo$, the connectivity ratio $\lamtwo/\lammax$, the
normalised spectral entropy of the eigenvalue distribution, the fraction of eigenvalue mass in the upper half of
the spectrum, and the spectral radius $\lammax$. Because all five come from the same decomposition, five cost
what one costs. The feature vector is the $L \times H \times 5$ array. We also evaluate each of the five alone,
to see which carries the signal.

\paragraph{Cross-layer dynamics (ours).}
Per-head differences between consecutive layers, $\phi_{\ell+1,h}-\phi_{\ell,h}$, for a chosen $\phi$. A flat
per-layer feature set represents depth only implicitly; the differences make the evolution explicit, and the
theory in Section~\ref{sec:theory:lapeig} predicts that a purely local feature family cannot express it.

\paragraph{Trajectory and velocity summaries.}
The head-averaged families that the prior literature uses, retained as ablations rather than as competitors:
per-layer metric values, trajectory statistics over layers (mean, dispersion, slope, range, extremal steps), and
layer-to-layer derivatives. Section~\ref{sec:exp:perhead} contrasts them with the per-head profile on identical
data.

\paragraph{LapEigvals \citep{binkowski2025lapeigvals}.}
The published attention-only state of the art, run from the authors' released implementation: per-layer,
per-head, position-normalised in-degree minus self-loop, sorted, top-$k$ kept, with $k$ selected on validation
from $\{5,10,25,50,100\}$ and a balanced logistic regression on the concatenation. Section~\ref{sec:theory:lapeig}
discusses what these features measure.

\subsection{Residual-stream detectors}
\label{sec:method:resid}

\paragraph{Token-role probe \citep{healy2026internal}.}
The published state of the art for this task. At eight evenly spaced depths, the residual state at the
function-name token, the mean over argument-value tokens, and the state at the closing delimiter are
concatenated, and a classifier is trained on the concatenation. We evaluate both the original single-hidden-layer
network and a regularised linear probe, and the token positions are located inside the generated span only.

\paragraph{Gram spectrum \citep{chen2024inside}.}
Eigenvalues of the covariance of the residual states over the generated span, together with effective rank and a
log-determinant surrogate, following the EigenScore family.

\subsection{Evaluation protocol}
\label{sec:method:protocol}

\paragraph{Grouped splits at the level of the tool.}
Examples are grouped by ground-truth tool name and split at the group level in a five-fold cross-fit, so every
score is on tools absent from training. This is stricter than splitting by prompt, and it is necessary: with
prompt-level splits, a classifier built from the tool identity alone reaches
\aucToolConfoundLlamaOneBGlaiveAll{} AUC on one of our runs, because errors concentrate on particular tools.
Under tool-level splits the same classifier is at or below chance everywhere.

\paragraph{Pooled cross-fit scoring.}
Each example is scored exactly once, by a model that never saw its tool, and we report one AUC pooled over the
whole dataset rather than the mean of per-fold AUCs, which on this data would be computed on folds containing a
handful of positives. We repeat the whole procedure over five split seeds and report the mean and standard
deviation across seeds.

\paragraph{Selection discipline.}
Every hyperparameter, the direction of every unsupervised score, and the $k$ of LapEigvals are chosen on
training and validation folds only. The one place a sweep over features appears, we select the feature on
training folds and report it on test, and label it as such.

\paragraph{Confound panel.}
Alongside the detectors we train three deliberately uninformative classifiers: prompt and generation lengths plus
a truncation flag, generation length alone, and a one-hot encoding of the ground-truth tool. A detector is only
interesting to the extent that it exceeds these.

\paragraph{Labels are recomputed at evaluation time.}
Labels are a function of the labelling code, not of the extraction run, so the evaluator re-derives every label
from the stored generation. This is what let us discover, after the fact, that an earlier labelling rule counted
a model's use of its own current date as a hallucination.
